% ============================================================
%  V. WHOLE-BODY RECONFIGURATION AND CONTACT CONTROL
% ============================================================
\section{Whole-Body Reconfiguration and Contact Control}
\label{sec:method}

The \ac{TERESA} system employs three complementary control strategies that
are active in different phases of the mission. During the approach, a
\ac{WBC} formulation drives the arm to maintain an active look-at
orientation toward the patient, feeding the anticipatory body scanner.
Between scanning targets, a pre-planned body reconfiguration step optimises
the legged base posture to bring each anatomical target into the arm's
optimal workspace. During scanning, a Cartesian impedance controller---built
specifically for the Z1 arm, which lacks native impedance hardware---maintains
compliant probe contact on the patient's skin. The handoff between these
controllers is managed by the mission \ac{FSM} described in
Section~\ref{sec:system_architecture}.

\subsection{Whole-Body Control for Reactive Arm Look-At}

During the approach phase, the Z1 arm must continuously orient its
\ac{EE}-mounted RealSense camera toward the patient's torso to feed the
anticipatory body scanner (Section~\ref{sec:anticipatory_scan}). The Spot
base simultaneously navigates toward the target fixed in the odometry frame
$\mathcal{F}_W$. \ac{TERESA} uses a \ac{WBC} formulation in which the arm
joint velocities are computed to track a desired \ac{EE} orientation, while
the base velocity is scaled by the perception-quality monitor.

\subsubsection{Arm look-at via damped pseudo-inverse}
Let $\mathbf{p}_{EE} \in \mathbb{R}^3$ be the current \ac{EE} position and
$\mathbf{p}^* \in \mathbb{R}^3$ the target position, both expressed in
$\mathcal{F}_W$. The desired \ac{EE} orientation is computed so that the
\ac{EE} $x$-axis points toward the target. Using the current arm
configuration $\mathbf{q}$ and the arm Jacobian
$\mathbf{J}_{\mathrm{arm}}(\mathbf{q}) \in \mathbb{R}^{6 \times 6}$
defined in~\eqref{eq:arm_jacobian}, the joint velocity command
$\dot{\mathbf{q}} \in \mathbb{R}^6$ is obtained via a damped
pseudo-inverse:
\begin{equation}
  \dot{\mathbf{q}} =
  \mathbf{J}_{\mathrm{arm}}^\top
  \!\left(
    \mathbf{J}_{\mathrm{arm}}\,\mathbf{J}_{\mathrm{arm}}^\top
    + \mu\,\mathbf{I}_6
  \right)^{-1}
  \mathbf{v}_{\mathrm{des}},
  \label{eq:damped_pinv}
\end{equation}
where $\mu > 0$ is a damping factor that ensures numerical stability near
kinematic singularities. The desired spatial velocity
$\mathbf{v}_{\mathrm{des}} \in \mathbb{R}^6$ is:
\begin{equation}
  \mathbf{v}_{\mathrm{des}} =
  \begin{bmatrix}
    k_r\,\mathbf{e}_{\mathrm{rot}} \\
    \mathbf{0}_3
  \end{bmatrix},
  \label{eq:vdes}
\end{equation}
where $k_r > 0$ is a rotational gain and $\mathbf{e}_{\mathrm{rot}} \in
\mathbb{R}^3$ is the orientation error between the current \ac{EE}
orientation and the desired look-at orientation. The orientation error is
computed using a minimum-rotation formulation: given the home \ac{EE}
orientation $\mathbf{R}_{\mathrm{home}} \in SO(3)$ and the desired look-at
direction $\mathbf{x}_{EE} = (\mathbf{p}^* - \mathbf{p}_{EE}) /
\|\mathbf{p}^* - \mathbf{p}_{EE}\|$, the desired rotation is the quaternion
that maps $\mathbf{R}_{\mathrm{home}}\,\mathbf{e}_x$ to $\mathbf{x}_{EE}$
with minimum angular displacement, relaxing the wrist configuration. The
translational component of $\mathbf{v}_{\mathrm{des}}$ is zero because the
translational target during scanning is provided separately by the body
scanner (Section~\ref{sec:anticipatory_scan}); during approach, the arm
maintains a nominal home position and contributes only orientation.

\subsubsection{Base velocity modulation}
The Spot base velocity, consisting of forward velocity $v_x$ and yaw rate
$\omega_z$, is controlled using a proportional law on the distance and
angle to the target in $\mathcal{F}_W$. These raw commands are then scaled
by the quality-aware velocity factor $v_{\mathrm{scale}}(q)$ derived from
the QualityMonitor in~\eqref{eq:velocity_scale}, as per
\eqref{eq:cmd_vel_scaled}. The base velocity is thus continuously modulated
by perceptual confidence, closing the whole-body perception-action loop.

\subsubsection{Robustness to TF degradation}
In a distributed system spanning the Spot onboard computer and an external
PC, the \ac{TF} tree can experience transient failures due to clock
desynchronisation or \ac{DDS} discovery delays. To maintain continuity of
base motion, the controller caches the last valid velocity command. If a
\ac{TF} lookup fails at the current control tick, the cached command is
re-published, preventing the base from stopping during brief connectivity
losses.

\subsection{Pre-Planned Body Reconfiguration for Scanning}

Once the anticipatory body scanner has localised the patient's anatomy and
computed the FAST target set $\mathcal{T}$ in $\mathcal{F}_W$
(Section~\ref{sec:anticipatory_scan}), the system must position the Z1 arm
to reach each target. Simply keeping Spot at a fixed posture for all
targets would force the arm to operate near the edge of its workspace for
non-central targets, degrading manipulability and positioning accuracy.
\ac{TERESA} addresses this by pre-planning an optimal body posture for
each scanning target before the arm moves.

\subsubsection{Workspace sweet spot}
The preferred arm configuration is defined by a \textit{sweet spot}
$\mathbf{p}_{\mathrm{sweet}} \in \mathbb{R}^3$ in the arm base frame
$\mathcal{F}_0$. This point represents the \ac{EE} position at which the
arm operates with high manipulability, comfortable joint angles, and
sufficient reserve to accommodate small patient movements. For the Z1 arm
mounted on Spot, $\mathbf{p}_{\mathrm{sweet}}$ is located approximately
in front of and slightly above the arm base.

\subsubsection{Grid search over body posture}
For each FAST target $\mathbf{p}_i^*$ expressed in the odometry frame
$\mathcal{F}_W$, the system searches over a discrete set of feasible body
postures. Let $\mathcal{H} = \{h_1, \ldots, h_{N_h}\}$ be a set of body
heights and $\mathcal{P} = \{\varphi_1, \ldots, \varphi_{N_p}\}$ a set of
pitch angles. For each candidate posture $(h, \varphi)$, the position of
the arm base frame $\mathcal{F}_0$ in $\mathcal{F}_W$ is computed from the
known static transform $\mathcal{F}_B \to \mathcal{F}_0$ and the current
body pose. The target position is then transformed to $\mathcal{F}_0$:
\begin{equation}
  \mathbf{p}_i^{\mathcal{F}_0}(h, \varphi) =
  \mathbf{T}_{\mathcal{F}_W}^{\mathcal{F}_0}(h, \varphi)\,
  \mathbf{p}_i^*.
  \label{eq:target_in_F0}
\end{equation}
The score for each posture is the negative Euclidean distance from the
transformed target to the sweet spot:
\begin{equation}
  s_i(h, \varphi) =
  -\bigl\|\,\mathbf{p}_i^{\mathcal{F}_0}(h, \varphi) -
  \mathbf{p}_{\mathrm{sweet}}\,\bigr\|.
  \label{eq:posture_score}
\end{equation}
The optimal posture for target $i$ is
$(h_i^*, \varphi_i^*) = \arg\max_{h \in \mathcal{H},\, \varphi \in \mathcal{P}}
s_i(h, \varphi)$, subject to the posture bounds
in~\eqref{eq:posture_bounds}. The yaw angle $\theta$ is never changed
during scanning, as it was already optimised during the approach phase.

\subsubsection{Inter-target execution}
The body reconfiguration is executed between scanning targets. Before the
arm moves to target $i$, Spot assumes the pre-computed posture
$(h_i^*, \varphi_i^*)$ and is given a settle interval
$\Delta t_{\mathrm{settle}}$ to allow the body to stabilise. The target
$\mathbf{p}_i^*$ is then re-transformed to $\mathcal{F}_0$ using the live
\ac{TF} tree, reflecting the actual posture achieved. The Z1 \ac{FSM}
receives a ``body ready'' signal and proceeds with the arm trajectory
toward the transformed target. After the arm completes the scan at target
$i$, the cycle repeats for target $i+1$. At the conclusion of the scan,
Spot returns to a nominal posture.

Because the entire posture optimisation is computed offline---a single grid
search over $\mathcal{H} \times \mathcal{P}$ for each of the $N$ targets,
performed immediately after the target set $\mathcal{T}$ is available---it
imposes no real-time computational burden and requires no online
optimisation during arm motion.

\subsection{Cartesian Impedance Control for Compliant Scanning}

The Unitree Z1 arm does not provide native impedance control: its
firmware supports only joint position tracking via the \ac{JTC} and raw
torque commands. To perform ultrasound scanning, however, the probe must
maintain gentle and continuous contact with the patient's skin, adapting
to surface curvature and small patient movements. \ac{TERESA} implements a
custom Cartesian impedance controller that runs directly on the Z1's
torque interface.

\subsubsection{Cartesian position control}
Let $\mathbf{p}_{EE}(t) \in \mathbb{R}^3$ be the current \ac{EE} position
and $\mathbf{p}_d(t) \in \mathbb{R}^3$ the desired position, updated at
\SI{500}{\hertz} from the RealSense surface estimate. The translational
position error is:
\begin{equation}
  \mathbf{e}_p(t) = \mathbf{p}_d(t) - \mathbf{p}_{EE}(t).
  \label{eq:pos_error}
\end{equation}
The Cartesian control force is computed as:
\begin{equation}
  \mathbf{F}_p =
  \mathbf{K}_p^{\mathrm{eff}}\,\mathbf{e}_p
  - \mathbf{K}_d^{\mathrm{eff}}\,\dot{\mathbf{p}}_{EE}
  + \mathbf{K}_i\,\int \mathbf{e}_p\,\mathrm{d}t,
  \label{eq:cartesian_force}
\end{equation}
where $\dot{\mathbf{p}}_{EE} = \mathbf{J}_{\mathrm{arm}}\,\dot{\mathbf{q}}$
is the \ac{EE} translational velocity, and
$\mathbf{K}_p^{\mathrm{eff}},\,\mathbf{K}_d^{\mathrm{eff}} \in
\mathbb{R}^{3 \times 3}$ are effective diagonal gain matrices that adapt
to the current arm extension:
\begin{equation}
  \mathbf{K}_p^{\mathrm{eff}} = (1 - 0.65\,\rho)\,\mathbf{K}_p,
  \qquad
  \mathbf{K}_d^{\mathrm{eff}} = (1 + \rho)\,\mathbf{K}_d,
  \label{eq:adaptive_gains}
\end{equation}
where $\rho = \mathrm{clip}\!\left(\|\mathbf{p}_{EE}^{xy}\|/r_{\max},
0, 1\right) \in [0,1]$ is the normalised planar reach, and
$r_{\max} = \SI{0.6}{\metre}$ is the arm's nominal planar reach. This
adaptive scaling reduces stiffness and increases damping as the arm
extends, improving stability near the workspace boundary. The nominal
gains are $\mathbf{K}_p = \mathrm{diag}(250, 250, 300)$,
$\mathbf{K}_d = \mathrm{diag}(15, 15, 15)$ in SI units, and
$\mathbf{K}_i = \mathrm{diag}(10, 10, 10)$.

\subsubsection{Surface-referenced desired position}
When the RealSense surface estimate is available, the desired position is
computed relative to the detected patient torso surface. Let
$\mathbf{p}_s \in \mathbb{R}^3$ be the surface centroid and
$\mathbf{n}_s \in \mathbb{S}^2$ the outward surface normal estimated
by the RealSense pipeline. The desired \ac{EE} position is:
\begin{equation}
  \mathbf{p}_d = \mathbf{p}_s + \delta_n\,\mathbf{n}_s,
  \label{eq:surface_target}
\end{equation}
where $\delta_n = \SI{-5}{\milli\metre}$ is a small offset along the
inward normal that ensures light probe contact without excessive force.

\subsubsection{Joint torque computation}
The Cartesian force $\mathbf{F}_p$ is projected into joint torques via
the arm Jacobian transpose, and combined with joint-space wrist control
and dynamic compensation:
\begin{equation}
  \boldsymbol{\tau} =
  \mathbf{J}_{\mathrm{arm}}^\top\,\tilde{\mathbf{F}}_p
  + \boldsymbol{\tau}_{\mathrm{wrist}}
  + \boldsymbol{\tau}_{\mathrm{comp}},
  \label{eq:torque_total}
\end{equation}
where $\tilde{\mathbf{F}}_p = [\mathbf{F}_p^\top,\,\mathbf{0}_3^\top]^\top
\in \mathbb{R}^6$ pads the translational force with zero moments.

The wrist torque $\boldsymbol{\tau}_{\mathrm{wrist}} \in \mathbb{R}^6$
applies joint-space \ac{PD} control to joints $q_4, q_5, q_6$ to maintain
the probe orientation perpendicular to the skin surface:
\begin{equation}
  \tau_{\mathrm{wrist},i} =
  K_{p,w}\,(q_{d,i} - q_i) - K_{d,w}\,\dot{q}_i,
  \quad i = 4,5,6,
  \label{eq:wrist_torque}
\end{equation}
with $K_{p,w} = \SI{40}{\newton\metre\per\radian}$ and
$K_{d,w} = \SI{4}{\newton\metre\second\per\radian}$.

The compensation torque $\boldsymbol{\tau}_{\mathrm{comp}} \in \mathbb{R}^6$
comprises gravity and Coriolis terms computed via the \acf{RNEA}
implemented in Pinocchio~\cite{carpentier2019}:
\begin{equation}
  \boldsymbol{\tau}_{\mathrm{comp}} =
  \boldsymbol{\tau}_{\mathrm{gravity}}(\mathbf{q})
  + \boldsymbol{\tau}_{\mathrm{Coriolis}}(\mathbf{q}, \dot{\mathbf{q}}).
  \label{eq:compensation}
\end{equation}

The total torque is saturated to the Z1 hardware limit
$\tau_{\max} = \SI{70}{\newton\metre}$ before being sent to the joint
controllers:
\begin{equation}
  \boldsymbol{\tau}_{\mathrm{cmd}} =
  \mathrm{clip}\!\left(
    \boldsymbol{\tau},\;
    -\tau_{\max}\,\mathbf{1},\;
    \tau_{\max}\,\mathbf{1}
  \right).
  \label{eq:torque_clip}
\end{equation}

\subsubsection{Controller switching}
The Z1 alternates between two \ac{ROS}~2 controllers: the \ac{JTC} for
position-controlled approach and homing motions, and the torque controller
for impedance-based scanning. A \texttt{safe\_controller\_switch} node
exposes two trigger services---\texttt{/safe\_switch/to\_torque} and
\texttt{/safe\_switch/to\_jtc}---that the \ac{FSM} calls before entering
and leaving the scanning phase. The \ac{JTC} is the safe default; the
system always returns to it on error. This custom impedance implementation
is essential to \ac{TERESA}'s scanning capability, as the Z1 arm's native
control stack provides no compliant contact mode.
